%%%%%%%%%%%%%%%%%%%%%%%%%%%%%%%%%%%%%%%%%%%%%%%%%%%%%%%%%%%%%%%%%%%%%%
%
% --------------------------------------------------------------------
%
% Document :  PDF test documentation document frame
%
% --------------------------------------------------------------------
%
%
%%%%%%%%%%%%%%%%%%%%%%%%%%%%%%%%%%%%%%%%%%%%%%%%%%%%%%%%%%%%%%%%%%%%%%

\nonstopmode

%%%%%%%%%%%%%%%% PDF %%%%%%%%%%%%%%%%%%%%%%%%%%%%%%
\documentclass[pdftex,12pt]{report}
\pdfoutput=1
\pdfimageresolution=288
\pdfcompresslevel=0

\usepackage{refmanVerifiedPDF}
\usepackage{verbatim}

\usepackage{url,xspace,boxedminipage}
\usepackage{graphicx,thumbpdf}
\usepackage[colorlinks,bookmarks,
  pdftitle={Test Results Documentation},
  pdfsubject={RT-Tester Test Documentation},
  pdfkeywords={},
  pdfauthor={Verified Systems}]{hyperref}

%%%%%%%%%%%%%%%%%%%%%%%%%%%%%%%%%%%%%%%%%%%%%%%%%%%

\usepackage[english]{babel}
\usepackage[latin1]{inputenc}
\usepackage[T1]{fontenc}
\usepackage{ae}
\usepackage{moreverb}
\usepackage{longtable}
\usepackage{lscape}
\usepackage{color}
%---  section numbering   ------------------------
\setcounter{tocdepth}{3}
\setcounter{secnumdepth}{3}

%-------------------------------------------------
% This is necessary to display �:
\DeclareInputText{181}{$\mu$}


%----------------------------------------------------

%%% \input{commands.tex}

%%% \definecolor{myblue}{rgb}{0.0,0.03,0.5}
\definecolor{myblue}{rgb}{0.0,0.0,1.0}

\newcommand{\com}[1]{\framebox
       {\parbox{\textwidth}{
           {\sf #1}
           }}}





\newcommand{\chg}[1]{\begin{tabular}{l|}
                      \parbox[t]{\textwidth}{ #1 } \\
                      \end{tabular}}

\newcommand{\percent}{\%}

\def\pict#1#2#3{
    \begin{figure}[htbp]
      \hfill\centerline{\includegraphics[width=\textwidth]{#2}}\hfill
      \caption{#1}\label{#3}
    \end{figure}
}


\def\pictsmall#1#2#3{
    \begin{figure}[htbp]
      \hfill\centerline{\includegraphics[width=5cm]{#2}}\hfill
      \caption{#1}\label{#3}
    \end{figure}
}

\def\pictmed#1#2#3{
    \begin{figure}[htbp]
      \hfill\centerline{\includegraphics[width=8cm]{#2}}\hfill
      \caption{#1}\label{#3}
    \end{figure}
}



% Verified's marginpar
\newcommand{\vmp}[1]{\marginpar{\raggedright{\textcolor{myblue}{\footnotesize\sf #1\normalsize}}}\ignorespaces}

\input{testrelatedflags.tex}


% 

\pagestyle{empty}

%----------------------------------------------------

\begin{document}
\begin{fullpage}
\pagestyle{empty}


\huge
{\bf Verified Systems International GmbH} \vskip -8mm
\rule{\textwidth}{0.5mm}
\vfill


\huge
{\bf 
     \_Project \\
     \\
     Test Results Document \\
     \ \\
     \aststlatitle\ (\aststla) \\
     \aststestname: \aststestheadline
     \ \\
     \astscidsbaseline

 }

\bigskip
\bigskip
\large
{\bf \makebox[7cm][l]{Test Execution Date:} \aststestexecutiondate}\\
{\bf \makebox[7cm][l]{Test Documentation Date:} \today}

\vfill
\footnotesize


Verified Systems International GmbH\\
Parkstra{\ss}e 123\\
28209 Bremen \\
Germany\\
\makebox[2cm][l]{Tel.  : } +49 421 57204-0 \\
\makebox[2cm][l]{Fax   : } +49 421 57204-22\\
\makebox[2cm][l]{e-mail: } info@verified.de\\
\makebox[2cm][l]{www: } http://www.verified.de

\makebox[7cm][l]{Author:} \astsauthor\\
\makebox[7cm][l]{Test execution performed by:} \aststestexecutionby\ on \aststestexecutiondate

\normalsize
\vfill
\vbox{\hsize=\textwidth{\rule{\textwidth}{0.3mm}} %\vskip 3pt
              \vskip -.6cm
              \hbox to \textwidth{\tiny \hfil 
              %%% \epsfxsize=1.8cm\epsfbox{pics/vlstamp_trans_red.eps}
               \includegraphics[width=1.8cm]{logo.pdf}
              }} 

\end{fullpage} 
%----------------------------------------------------
\newpage
\pagestyle{myheadings}
\markright{RT-Tester 6.0 TQ -- \aststla\ \aststestname\ --  Test Results Documentation -- \astscidsbaseline}

\tableofcontents
%%% \listoffigures


%------------------------------------------------------------------
\input{testobj.tex}
%-------------------------------------------------------

\chapter{Test Results}

%...................................................................
\section{Discrepancies}\label{sec:discrep}

\input{discrepancies.tex}

%...................................................................
\input{remarks.tex}

%...................................................................
\section{Result from Automatic Test Execution}


\small
\input{GENERATEDVERDICT.tex}
\normalsize
\bigskip

%...................................................................
\section{Overall Test Result}

The verdict established by the automatic test execution may be overridden due
to manual analysis. If this is the case, the reasons are given in
Section~\ref{sec:discrep}. 

\bigskip

\input{verdict.tex}

%...................................................................
\chapter{Test Case Coverage}

\input{coverage.tex}


%-------------------------------------------------------------------
\begin{fullpage} 
\chapter{How to Read the Test Execution Log}

The test execution log provides near-complete information in a single file.
Every entry is preceded by a line of the following form.
\begin{verbatim}
TM 00000000000 UNIT O 
\end{verbatim}
The number after TM is the time stamp, counted in seconds. 
%% @todo: update when time tick is im milli/micro
During unit tests, only output ("O") entries exist.

\section{Test Context Information}

The startup information of a test informs about the
user, application, and test case.
The local time of the test start is documented in the following form.
\begin{verbatim}
START OF TEST : <YEAR>:<MONTH>:<DAY>:<HOUR>:<MINUTE>:<SECONDS>
\end{verbatim}

The test termination information summarises
the number of
errors   (failures of @rttAssert conditions),
and the number of 
warnings (failures of @rttCheck conditions).
The verdict of the test run has one of the following values.

\begin{tabular}{l@{~:~~~}l}
{\tt PASS}       & successful execution of the test; no errors; potentially warnings 
\\
{\tt FAIL}       & test execution did yield one or more errors
\\
{\tt INCONC}     & inconclusive: test run was not finished properly
\\
{\tt TESTERROR}  & test execution failed, the test run is unreliable
\\
\end{tabular}

In case of {\tt TESTERROR} the observed information on the SUT cannot be used.
Some typical scenarios for this situations are
\begin{itemize}
\item   interrupt signals were caught,
\item   files could not be opened, 
\item  a user-defined
   {\tt @rtt\-Test\-Error} 
   situation was reached
\end{itemize}

%% -------------------------------------------------------------

\section{Test Step Information}

\begin{itemize}
\item start of a test step, indicated by a block of the form
\begin{verbatim}
============================================================================
======================== BEGIN of RTT Test Step <N> ========================
============================================================================
\end{verbatim}
\item formatted entries, where the following information may be 
  provided by the user
  \begin{itemize}
  \item {\tt CONDITION} : precondition        (generated situation of the SUT)
  \item {\tt EVENT}     : execution situation (e.g, a function call)
  \item {\tt EXPECTED}  : expected SUT behaviour
  \item {\tt TAG}       : test case tag this test step corresponds to
  \end{itemize}
\item comparison of expected and observed behaviour (@rttAssert)
\item end of a test step, indicated by a block of the form
\begin{verbatim}
------------------------  END  of RTT Test Step <N> ------------------------
\end{verbatim}
\end{itemize}

%% -------------------------------------------------------------

\section{Stub Call Information}

Calls to stub functions are documented in the test execution log in 
the following form:
\begin{verbatim}
TM 00000000000 UNIT O CALL of stub function <function name>(...):
                       <QUALIFIER> <variable>     = <value>
                      --------------------------------
                       <QUALIFIER> <variable>     = <value>
                      ========================================================
\end{verbatim}
The line ``{\tt ------------}'' separates input and output parameters.
The {\tt \verb+<+QUALIFIER\verb+>+} information is one of 
{\tt IN}, {\tt OUT}, or {\tt INOUT},  or {\tt RETURN}.
The {\tt IN}, {\tt OUT}, {\tt INOUT} information has to be provided in the stub declaration 
(see the {\tt .\-stubs} file), {\tt RETURN} denotes the return value of the
 parameter. 


For pointer values, by convention the variable names are preceded by a
``{\tt *}'' and the values pointed to are printed.
E.g, for a function {\tt a\_\-pointer} with the stub signature
\begin{verbatim}
  int* a_pointer(@inout int* counter_p);
\end{verbatim}
the following log entry may be generated.
\begin{verbatim}
TM 00000000000 UNIT O CALL of stub function a_pointer(...):
                       INOUT *counter_p =    17
                      --------------------------------
                       INOUT *counter_p =    18
                      RETURN *a_pointer =     2
                      ========================================================
\end{verbatim}


%----------------------------------------------------------------------------
\input{glossary.tex}

%----------------------------------------------------------------------------
\begin{landscape}
\chapter{Complete Test Execution Log}

\footnotesize
\verbatiminput{\aststestexecutionlog}
\normalsize
\end{landscape}
%-----------------------------------------------------------------
\input{astsversion.tex}

%-----------------------------------------------------------------
\input{testreportappendix.tex}

\end{fullpage} 
%-----------------------------------------------------------------
\end{document}

%% -------------------------------------------------------------------------
% LocalWords:  SUT rttAssert rttCheck TM CIDS GmbH
